\small
\begin{enumerate}[label=\textbf{\arabic*.}, leftmargin=0pt, itemsep=1.2em]
\item Что означает аббревиатура ADALINE?
\begin{itemize}[label=$\square$, leftmargin=1.5cm, itemsep=0.2em, topsep=0.3em]
\item Advanced Data Layer Intelligence
\item Adaptive Linear Neuron
\item Artificial Data Logical Network
\item Automatic Linear Equation
\end{itemize}
{\small\textit{Правильный ответ: Adaptive Linear Neuron.}}

\item Кем была разработана архитектура сети ADALINE?
\begin{itemize}[label=$\square$, leftmargin=1.5cm, itemsep=0.2em, topsep=0.3em]
\item Марвином Минским
\item Бернардом Уидроу и Марцианом Хоффом
\item Фрэнком Розенблаттом
\item Джоном Хопфилдом
\end{itemize}
{\small\textit{Правильный ответ: Бернардом Уидроу и Марцианом Хоффом.}}

\item В чем состоит главное отличие ADALINE от классического персептрона Розенблатта в процессе обучения?
\begin{itemize}[label=$\square$, leftmargin=1.5cm, itemsep=0.2em, topsep=0.3em]
\item ADALINE использует нелинейную функцию активации
\item В ADALINE ошибка вычисляется до применения пороговой функции (по линейному выходу)
\item ADALINE не имеет весов
\item ADALINE обучается без учителя
\end{itemize}
{\small\textit{Правильный ответ: В ADALINE ошибка вычисляется до применения пороговой функции (по линейному выходу).}}

\item Какая функция активации используется в ADALINE на этапе вычисления ошибки для подстройки весов?
\begin{itemize}[label=$\square$, leftmargin=1.5cm, itemsep=0.2em, topsep=0.3em]
\item Сигмоидальная
\item Единичная ступенчатая
\item Линейная
\item Релу (ReLU)
\end{itemize}
{\small\textit{Правильный ответ: Линейная.}}

\item Какую функцию потерь (cost function) обычно минимизируют при обучении ADALINE?
\begin{itemize}[label=$\square$, leftmargin=1.5cm, itemsep=0.2em, topsep=0.3em]
\item Кросс-энтропию (Cross-entropy)
\item Сумму или среднее квадратов ошибок (MSE)
\item Среднюю абсолютную ошибку (MAE)
\item Коэффициент детерминации
\end{itemize}
{\small\textit{Правильный ответ: Сумму или среднее квадратов ошибок (MSE).}}

\item Как называется правило обучения, применяемое в ADALINE?
\begin{itemize}[label=$\square$, leftmargin=1.5cm, itemsep=0.2em, topsep=0.3em]
\item Правило Хебба
\item Алгоритм обратного распространения ошибки
\item Правило Уидроу-Хоффа (Дельта-правило)
\item Алгоритм Кохонена
\end{itemize}
{\small\textit{Правильный ответ: Правило Уидроу-Хоффа (Дельта-правило).}}

\item Каким методом осуществляется минимизация функции потерь в итеративном алгоритме ADALINE?
\begin{itemize}[label=$\square$, leftmargin=1.5cm, itemsep=0.2em, topsep=0.3em]
\item Методом главных компонент
\item Алгоритмом k-средних
\item Градиентным спуском (Gradient Descent)
\item Генетическим алгоритмом
\end{itemize}
{\small\textit{Правильный ответ: Градиентным спуском (Gradient Descent).}}

\item Какое свойство функции квадратичной ошибки делает градиентный спуск эффективным для ADALINE?
\begin{itemize}[label=$\square$, leftmargin=1.5cm, itemsep=0.2em, topsep=0.3em]
\item Функция является выпуклой и имеет только один глобальный минимум
\item Функция имеет множество локальных минимумов
\item Функция недифференцируема
\item Функция никогда не достигает нуля
\end{itemize}
{\small\textit{Правильный ответ: Функция является выпуклой и имеет только один глобальный минимум.}}

\item В контексте прогнозирования временных рядов, что такое «размер окна» (window size)?
\begin{itemize}[label=$\square$, leftmargin=1.5cm, itemsep=0.2em, topsep=0.3em]
\item Количество нейронов в скрытом слое
\item Количество предыдущих отсчетов времени, используемых для предсказания следующего значения
\item Количество эпох обучения
\item Скорость обучения
\end{itemize}
{\small\textit{Правильный ответ: Количество предыдущих отсчетов времени, используемых для предсказания следующего значения.}}

\item Как формируется обучающая выборка для прогнозирования временного ряда с помощью ADALINE?
\begin{itemize}[label=$\square$, leftmargin=1.5cm, itemsep=0.2em, topsep=0.3em]
\item Данные случайно перемешиваются без учета времени
\item Применяется метод скользящего окна, где входы – это история ряда, а выход – следующее значение
\item Используется только первое значение ряда для предсказания всех остальных
\item Весь ряд подается на один вход нейрона
\end{itemize}
{\small\textit{Правильный ответ: Применяется метод скользящего окна, где входы – это история ряда, а выход – следующее значение.}}

\item Помимо итеративного (градиентного) метода, как еще можно найти оптимальные веса для линейной нейронной сети?
\begin{itemize}[label=$\square$, leftmargin=1.5cm, itemsep=0.2em, topsep=0.3em]
\item Только полным перебором
\item Аналитически, используя псевдообратную матрицу
\item С помощью дерева решений
\item Путем случайной генерации
\end{itemize}
{\small\textit{Правильный ответ: Аналитически, используя псевдообратную матрицу.}}

\item Что такое псевдообратная матрица (матрица Мура-Пенроуза)?
\begin{itemize}[label=$\square$, leftmargin=1.5cm, itemsep=0.2em, topsep=0.3em]
\item Матрица, состоящая только из нулей
\item Обобщение понятия обратной матрицы для неквадратных матриц
\item Матрица случайных чисел
\item Единичная матрица отрицательного размера
\end{itemize}
{\small\textit{Правильный ответ: Обобщение понятия обратной матрицы для неквадратных матриц.}}

\item Какое выражение используется для аналитического вычисления оптимальных весов линейной нейронной сети через псевдообратную матрицу?
\begin{itemize}[label=$\square$, leftmargin=1.5cm, itemsep=0.2em, topsep=0.3em]
\item Произведение псевдообратной матрицы входных данных на вектор целевых значений
\item Простое произведение матрицы входных данных на вектор целевых значений
\item Сумма псевдообратной матрицы и вектора целевых значений
\item Произведение транспонированного вектора целевых значений на матрицу входных данных
\end{itemize}
{\small\textit{Правильный ответ: Произведение псевдообратной матрицы входных данных на вектор целевых значений.}}

\item В чем преимущество аналитического решения (через псевдообратную матрицу) перед градиентным спуском?
\begin{itemize}[label=$\square$, leftmargin=1.5cm, itemsep=0.2em, topsep=0.3em]
\item Оно требует меньше оперативной памяти для огромных датасетов
\item Оно находит точное решение за один шаг (без подбора learning rate и эпох)
\item Оно позволяет использовать сложные нелинейные функции активации
\item Оно не требует численных библиотек
\end{itemize}
{\small\textit{Правильный ответ: Оно находит точное решение за один шаг (без подбора learning rate и эпох).}}

\item Какое ключевое свойство псевдообратной матрицы позволяет использовать её для нахождения оптимальных весов линейной нейронной сети?
\begin{itemize}[label=$\square$, leftmargin=1.5cm, itemsep=0.2em, topsep=0.3em]
\item Она гарантирует нулевую ошибку для любых данных
\item Она даёт решение по методу наименьших квадратов, минимизируя сумму квадратов отклонений
\item Она работает только с квадратными матрицами
\item Она требует обратимости матрицы признаков
\end{itemize}
{\small\textit{Правильный ответ: Она даёт решение по методу наименьших квадратов, минимизируя сумму квадратов отклонений.}}

\item Если обучающая выборка содержит мультиколлинеарные признаки (сильно зависимые), какой метод вычисления весов будет более стабильным?
\begin{itemize}[label=$\square$, leftmargin=1.5cm, itemsep=0.2em, topsep=0.3em]
\item Обычная обратная матрица (inv)
\item Псевдообратная матрица (pinv), основанная на сингулярном разложении (SVD)
\item Использование случайных весов
\item Отбрасывание всех признаков
\end{itemize}
{\small\textit{Правильный ответ: Псевдообратная матрица (pinv), основанная на сингулярном разложении (SVD).}}

\item Как вычисляется среднеквадратичная ошибка (MSE)?
\begin{itemize}[label=$\square$, leftmargin=1.5cm, itemsep=0.2em, topsep=0.3em]
\item Как среднее суммы модулей разностей прогноза и факта
\item Как корень из максимальной ошибки
\item Как среднее значение квадратов разностей между прогнозируемым и фактическим значениями
\item Как процент правильных классификаций
\end{itemize}
{\small\textit{Правильный ответ: Как среднее значение квадратов разностей между прогнозируемым и фактическим значениями.}}

\item Чем отличается прогнозирование ADALINE от задачи классификации персептроном?
\begin{itemize}[label=$\square$, leftmargin=1.5cm, itemsep=0.2em, topsep=0.3em]
\item ADALINE выдает непрерывное число (прогноз), а персептрон – дискретный класс (например, 0 или 1)
\item ADALINE обучается дольше персептрона
\item Персептрон прогнозирует непрерывные значения
\item Отличий нет
\end{itemize}
{\small\textit{Правильный ответ: ADALINE выдает непрерывное число (прогноз), а персептрон – дискретный класс (например, 0 или 1).}}

\item Если скорость обучения в градиентном спуске слишком велика, что произойдет с ошибкой?
\begin{itemize}[label=$\square$, leftmargin=1.5cm, itemsep=0.2em, topsep=0.3em]
\item Ошибка будет плавно уменьшаться
\item Ошибка может начать расходиться (стремиться к бесконечности)
\item Ошибка сразу станет нулевой
\item Никаких изменений не будет
\end{itemize}
{\small\textit{Правильный ответ: Ошибка может начать расходиться (стремиться к бесконечности).}}

\item Как влияет нестандартизованность/немасштабированность признаков на градиентный спуск в ADALINE?
\begin{itemize}[label=$\square$, leftmargin=1.5cm, itemsep=0.2em, topsep=0.3em]
\item Увеличивает скорость сходимости
\item Делает поверхность ошибки вытянутой, что замедляет сходимость алгоритма
\item Никак не влияет
\item Гарантирует нулевую ошибку
\end{itemize}
{\small\textit{Правильный ответ: Делает поверхность ошибки вытянутой, что замедляет сходимость алгоритма.}}

\item Что означает аббревиатура SVD в контексте вычисления псевдообратной матрицы?
\begin{itemize}[label=$\square$, leftmargin=1.5cm, itemsep=0.2em, topsep=0.3em]
\item Standard Vector Derivative
\item Singular Value Decomposition (Сингулярное разложение)
\item Super Vector Data
\item Support Vector Design
\end{itemize}
{\small\textit{Правильный ответ: Singular Value Decomposition (Сингулярное разложение).}}

\item Какова геометрическая интерпретация выхода ADALINE с двумя входами?
\begin{itemize}[label=$\square$, leftmargin=1.5cm, itemsep=0.2em, topsep=0.3em]
\item Уравнение параболы
\item Уравнение плоскости в трехмерном пространстве
\item Сфера
\item Точка
\end{itemize}
{\small\textit{Правильный ответ: Уравнение плоскости в трехмерном пространстве.}}

\item Как добавить возможность ADALINE учитывать смещение (bias) при матричных вычислениях?
\begin{itemize}[label=$\square$, leftmargin=1.5cm, itemsep=0.2em, topsep=0.3em]
\item Вычесть единицу из всех ответов
\item Добавить к матрице признаков X столбец, состоящий из единиц
\item Умножить матрицу X на 0
\item Это невозможно сделать
\end{itemize}
{\small\textit{Правильный ответ: Добавить к матрице признаков X столбец, состоящий из единиц.}}

\item При прогнозировании на несколько шагов вперед рекуррентным методом (заменяя неизвестные значения собственными прогнозами), как ведет себя ошибка со временем?
\begin{itemize}[label=$\square$, leftmargin=1.5cm, itemsep=0.2em, topsep=0.3em]
\item Ошибка остается постоянной
\item Ошибка уменьшается
\item Ошибка накапливается (экспоненциально возрастает)
\item Ошибка становится отрицательной
\end{itemize}
{\small\textit{Правильный ответ: Ошибка накапливается (экспоненциально возрастает).}}

\item Какой метод обновления весов обновляет веса после предъявления каждого отдельного примера из выборки?
\begin{itemize}[label=$\square$, leftmargin=1.5cm, itemsep=0.2em, topsep=0.3em]
\item Пакетный градиентный спуск (Batch Gradient Descent)
\item Стохастический градиентный спуск (Stochastic Gradient Descent)
\item Вычисление аналитического решения
\item Генетический алгоритм
\end{itemize}
{\small\textit{Правильный ответ: Стохастический градиентный спуск (Stochastic Gradient Descent).}}

\item Какая метрика покажет, на сколько (например, в рублях или градусах) ошибается наша линейная модель в среднем без возведения в квадрат?
\begin{itemize}[label=$\square$, leftmargin=1.5cm, itemsep=0.2em, topsep=0.3em]
\item MSE
\item MAE (Mean Absolute Error)
\item R\\textasciicircum{}2
\item Кросс-энтропия
\end{itemize}
{\small\textit{Правильный ответ: MAE (Mean Absolute Error).}}

\item Когда аналитическое вычисление весов становится вычислительно неэффективным?
\begin{itemize}[label=$\square$, leftmargin=1.5cm, itemsep=0.2em, topsep=0.3em]
\item Когда число признаков очень мало
\item Когда число примеров мало
\item Когда матрица признаков имеет огромную размерность (например, миллион признаков), так как вычислительная сложность обращения матрицы растёт кубически с размерностью
\item Когда данные линейно разделимы
\end{itemize}
{\small\textit{Правильный ответ: Когда матрица признаков имеет огромную размерность (например, миллион признаков), так как вычислительная сложность обращения матрицы растёт кубически с размерностью.}}

\item Может ли ADALINE идеально прогнозировать сложный нелинейный временной ряд (например, синусоиду) без добавления нелинейных признаков?
\begin{itemize}[label=$\square$, leftmargin=1.5cm, itemsep=0.2em, topsep=0.3em]
\item Да, так как работает с плавающей точкой
\item Нет, она может построить только линейный тренд
\item Да, если обучать очень долго
\item Да, если использовать псевдообратную матрицу
\end{itemize}
{\small\textit{Правильный ответ: Нет, она может построить только линейный тренд.}}

\item В чем отличие дельта-правила от правила обучения персептрона?
\begin{itemize}[label=$\square$, leftmargin=1.5cm, itemsep=0.2em, topsep=0.3em]
\item Дельта-правило всегда использует целочисленные математические операции
\item Дельта-правило обновляет веса пропорционально непрерывной ошибке, а классический персептрон – только при смене дискретного класса
\item Никакой разницы нет
\item Дельта-правило нельзя применять программно
\end{itemize}
{\small\textit{Правильный ответ: Дельта-правило обновляет веса пропорционально непрерывной ошибке, а классический персептрон – только при смене дискретного класса.}}

\item Что означает отрицательное значение веса при прогнозировании временного ряда с помощью ADALINE?
\begin{itemize}[label=$\square$, leftmargin=1.5cm, itemsep=0.2em, topsep=0.3em]
\item Это ошибка компилятора
\item Данный предыдущий отсчет имеет обратную (отрицательную) корреляцию с прогнозируемым значением
\item Вес не влияет на прогноз
\item Ряд не имеет решения
\end{itemize}
{\small\textit{Правильный ответ: Данный предыдущий отсчет имеет обратную (отрицательную) корреляцию с прогнозируемым значением.}}

\end{enumerate}
